\section{Read a real timing result without changing its pixels}
The deterministic clock makes exact arithmetic observable. An ordinary CLI run
instead uses the host clock and actual sleep service. The following observations
come from one fresh process for each input and start position on the author's
machine. The caller completed the actual encode, store, reload, validation and
playback path. Final image dumps matched the corresponding source frames.

\begin{table}[htbp]
\centering
\caption{Reported playback state from four real-clock CLI runs; times are milliseconds.}
\label{tab:video-playback-host-observation}
\begin{tabular}{l r r r r r}
\toprule
Input & Start & Frames & Logical & Late intervals & Maximum late \\
\midrule
moving & 0 & 8 & 320 & 8 & 45 \\
moving & 6 & 2 & 80 & 2 & 27 \\
still & 0 & 4 & 160 & 4 & 23 \\
all change & 1 & 1 & 40 & 1 & 20 \\
\bottomrule
\end{tabular}
\end{table}

Each run retains the logical duration implied by its frame suffix. None reports
an exact host schedule. An initial seek can finish late, and a requested sleep
can overshoot; lateness is recorded rather than hidden by moving the anchor.
A late interval is still an accepted frame with the required numeric pixels.
Do not use these particular lateness counts as expected answers on another run.
The maximum values are observations of these runs, not bounds on the simulator.

Loading and initial validation happen before the anchor; the first seek and later
record applications happen inside the schedule. Whole-process elapsed time also
includes JVM startup, assembly, encoding and file I/O, so it cannot stand in for
isolated playback cost. No IntelliJ click-to-frame or repaint time was measured
here. Figure~\ref{fig:video-playback-progress} separately records an actual IDE
publication state; it does not measure click-to-frame or repaint latency.
